\documentclass[11pt]{article}
\usepackage[margin=1in]{geometry}
\usepackage{amsmath,amssymb,amsthm}
\usepackage{booktabs}
\usepackage[hidelinks]{hyperref}

\newtheorem{theorem}{Theorem}[section]
\newtheorem{proposition}[theorem]{Proposition}
\newtheorem{lemma}[theorem]{Lemma}
\newtheorem{corollary}[theorem]{Corollary}
\theoremstyle{remark}
\newtheorem{remark}[theorem]{Remark}

\newcommand{\Z}{\mathbb{Z}}
\newcommand{\F}{\mathbb{F}}

\title{Towards a direct construction of a 5-GDD of type $10^{27}$:\\
obstructions to symmetric constructions and a computer search}
\author{}
\date{}

\begin{document}
\maketitle

\begin{abstract}
The existence of a group divisible design with block size five and type $10^{27}$
is one of the cases left open by Forbes~\cite{Forbes2023}.
We attempted a direct (computer) construction.
\textbf{We did not find such a design}; this note records what we can prove and what we tried.
We show that no 5-GDD of type $10^{27}$ admits an automorphism group acting regularly on
its $270$ points (in particular none is cyclic), that none admits the group
$\{x\mapsto qx+t\}$ of $\F_{27}$ ($q$ a nonzero square) acting on the groups, and that
in the natural framework used by Forbes for type $10^{15}$
(point set $\Z_{270}$, groups the residue classes modulo $27$, development
$x\mapsto x+2j$) no further affine automorphism can be imposed, so that $26$ base blocks
are unavoidable there.  We also derive a necessary ``signed difference'' identity for that
framework and summarise the computer searches that were carried out.
\end{abstract}

\section{Introduction}

A \emph{$5$-GDD of type $g^u$} is a triple $(V,\mathcal G,\mathcal B)$ where $V$ is a set of
$gu$ points, $\mathcal G$ is a partition of $V$ into $u$ \emph{groups} of size $g$, and
$\mathcal B$ is a collection of $5$-subsets of $V$ (\emph{blocks}) such that every pair of
points from distinct groups lies in exactly one block and no pair from the same group lies
in a block.  The necessary conditions
\[
u\ge 5,\qquad g(u-1)\equiv 0 \pmod 4,\qquad g^2u(u-1)\equiv 0 \pmod{20}
\]
hold for $(g,u)=(10,27)$, and by~\cite[Theorem~1.2]{Forbes2023} the case $10^{27}$ is one
of the few remaining possible exceptions for $g=10$ (together with $u\in\{5,7,39,47\}$).
A 5-GDD of type $10^{27}$ has
\[
v=270,\qquad b=\frac{270\cdot 260}{2\cdot 10}=3510,\qquad r=\frac{260}{4}=65 .
\]
If it existed, Lemma~2.1 of~\cite{Forbes2023} (inflation by a factor $5$) would also give
type $50^{27}$, the remaining exception for $g=50$.

\medskip\noindent
\textbf{Status.} No 5-GDD of type $10^{27}$ is constructed in this note.
Sections~\ref{sec:regular}--\ref{sec:F27} contain complete proofs of non-existence of designs
with certain prescribed automorphism groups; Section~\ref{sec:search} describes the
computer search and its (negative) outcome.

\section{No regular automorphism group}\label{sec:regular}

\begin{proposition}\label{prop:regular}
No $5$-GDD of type $10^{27}$ admits an automorphism group acting regularly on its $270$ points.
In particular there is no cyclic $5$-GDD of type $10^{27}$ (one invariant under $x\mapsto x+1$ on $\Z_{270}$).
\end{proposition}

\begin{proof}
Suppose $G$, $|G|=270$, acts regularly; identify $V$ with $G$ (acting by left multiplication).
The group containing the identity is a subgroup $H$ of order $10$ and the groups are the left
cosets $gH$.  An element of order $2$ acts on $G$ as a product of $135$ transpositions, an odd
permutation, so the even permutations in $G$ form a subgroup $G_0$ of index $2$, and the
involution of $H$ lies outside $G_0$; hence every coset $gH$ meets $G_0$ in exactly $5$ points.

Call a pair of points \emph{split} if exactly one of them lies in $G_0$.  The number of split
pairs lying in distinct groups is
\[
135\cdot 135-27\cdot 5\cdot 5=17550 .
\]
A block $B$ with $e$ points in $G_0$ and $o=5-e$ points outside contains $eo\in\{0,4,6\}$ split
pairs, and $eo$ is constant on the $G$-orbit of $B$ (elements of $G_0$ preserve $G_0$, the others
interchange $G_0$ and its complement).  The stabiliser of a block acts semiregularly on it, so it
has order $1$ or $5$.  In the latter case the block is an orbit of a subgroup of order $5$, which
lies in $G_0$, so the block lies in one coset of $G_0$ and $eo=0$.  Hence, summing over
representatives $B_1,\dots,B_s$ of the full orbits,
\[
17550=270\sum_{i=1}^{s} e_io_i ,\qquad\text{i.e.}\qquad \sum_{i=1}^{s}e_io_i=65,
\]
which is impossible because every $e_io_i$ is even.
\end{proof}

This explains why the construction of type $10^{15}$ in~\cite{Forbes2023} uses the mapping
$x\mapsto x+2j$ rather than $x\mapsto x+j$; the same parity argument applies to
type $10^{u}$ whenever $u\equiv 3\pmod 4$.

\section{\texorpdfstring{The framework $\Z_{270}$, $x\mapsto x+2j$}{The framework Z270, x -> x+2j}}\label{sec:framework}

Let $V=\Z_{270}$ with groups the residue classes modulo $27$, and look for a design invariant under
$\tau\colon x\mapsto x+2$.  By the Chinese remainder theorem we write $x\in\Z_{270}$ as
$(a,\ell)\in\Z_{135}\times\{0,1\}$ with $a=x \bmod 135$ and $\ell=x\bmod 2$ (the \emph{level});
then $\tau$ acts as $(a,\ell)\mapsto(a+2,\ell)$, i.e.\ $\langle\tau\rangle\cong\Z_{135}$ acts
regularly on each level, and two points lie in the same group if and only if their
$a$-coordinates agree modulo $27$.  Put $H=27\Z_{135}$.

No block is fixed by a nontrivial element of $\langle\tau\rangle$ (a fixed block would be a
union of orbits of a subgroup of order $3$, or of $H$, which is impossible for $5$ points from
distinct groups), so the design consists of
\[
3510/135=26
\]
full orbits.  For a base block $B$ write $B^0,B^1\subseteq\Z_{135}$ for its level-$0$ and level-$1$
coordinates and $e_B=|B^0|$, $o_B=|B^1|$.

\begin{lemma}\label{lem:diff}
Base blocks $B_1,\dots,B_{26}$ generate a $5$-GDD of type $10^{27}$ under $\langle\tau\rangle$ if and only if
\begin{enumerate}
\item the differences $x-y$ ($x\ne y\in B_i^0$, all $i$) are exactly the elements of $\Z_{135}\setminus H$, each once;
\item the same holds for the level-$1$ parts $B_i^1$;
\item the \emph{mixed} differences $x-y$ ($x\in B_i^1$, $y\in B_i^0$, all $i$) are exactly the elements of $\Z_{135}\setminus H$, each once.
\end{enumerate}
Consequently $\sum e_B(e_B-1)=\sum o_B(o_B-1)=\sum e_Bo_B=130$, hence $\sum e_B=\sum o_B=65$ and
$\sum e_B^2=195$.
\end{lemma}

\begin{proof}
Standard: a pair $\{(a,\ell),(a',\ell')\}$ is covered once by the developed design exactly when the
corresponding difference occurs once in the appropriate list, and pairs inside a group correspond
to differences in $H$.  The counts follow from $|\Z_{135}\setminus H|=130$.
\end{proof}

For example, the patterns $(n_{4,1},n_{3,2},n_{2,3},n_{1,4})=(6,8,5,7)$ or $(7,5,8,6)$
(numbers of base blocks with $(e,o)$ points on the two levels) satisfy the counting conditions.

\begin{proposition}[Signed difference identity]\label{prop:signed}
For a base block $B$ let $D_B=\sum_{x\in B^0}[\bar x]-\sum_{y\in B^1}[\bar y]\in\Z[\Z_{27}]$, where
$\bar x=x\bmod 27$.  Then
\[
\sum_{i=1}^{26} D_{B_i}D_{B_i}^{(-1)}=130 .
\]
Equivalently, the $27\times 702$ matrix with rows indexed by $r\in\Z_{27}$, columns by pairs
$(i,x)\in\{1,\dots,26\}\times\Z_{27}$ and entries the coefficient of $[x-r]$ in $D_{B_i}$, has
pairwise orthogonal rows of weight $130$.
\end{proposition}

\begin{proof}
Reduce Lemma~\ref{lem:diff} modulo $27$: each nonzero residue occurs exactly $5$ times among the level-$0$
differences, $5$ times among the level-$1$ differences and $5$ times among the mixed differences, and $0$
occurs only as $x-x$.  Writing $A_i,C_i$ for the level parts,
$\sum A_iA_i^{(-1)}=65+5(\Z_{27}-1)$, $\sum C_iC_i^{(-1)}=65+5(\Z_{27}-1)$ and
$\sum C_iA_i^{(-1)}=\sum A_iC_i^{(-1)}=5(\Z_{27}-1)$, where $\Z_{27}$ also denotes the sum of all group
elements.  Hence $\sum D_iD_i^{(-1)}=130$.
\end{proof}

We checked that the published design of type $10^{15}$~\cite{Forbes2023}, which has the same form
with $\Z_{150}$ and residue classes modulo $15$, satisfies the analogous identity
$\sum D_BD_B^{(-1)}=70$ in $\Z[\Z_{15}]$.

\section{No further affine automorphisms in the framework}\label{sec:affine}

\begin{proposition}\label{prop:affine}
Let $\mathcal D$ be a $5$-GDD of type $10^{27}$ on $\Z_{270}$ with groups the residue classes modulo $27$,
invariant under $x\mapsto x+2$.  If $\varphi(x)=mx+t$ with $\gcd(m,270)=1$ is an automorphism of
$\mathcal D$, then $m=1$ and $t$ is even.
\end{proposition}

\noindent
Thus the base blocks cannot be organised into orbits of a larger affine group, and $26$ base blocks
are unavoidable in this framework.  We first record a simple observation.

\begin{lemma}\label{lem:transl}
Let $\mu$ be a level-preserving automorphism of the form $(a,\ell)\mapsto(ma,\ell)$, let $B$ be a block and
let $P,Q\in B$ lie on the same level.  If $\mu(Q)-\mu(P)=\pm(Q-P)$, or if $P,Q$ lie on different levels and
$\mu(Q)-\mu(P)=Q-P$ (differences of $a$-coordinates), then $\mu(B)=B+s$ for some $s\in\Z_{135}$.
\end{lemma}

\begin{proof}
In each case $\{\mu(P),\mu(Q)\}$ is a translate $\{P,Q\}+s$ of $\{P,Q\}$ by an element of
$\Z_{135}$.  The blocks $\mu(B)$ and $B+s$ both contain this pair, so they coincide.
\end{proof}

\begin{proof}[Proof of Proposition~\ref{prop:affine}]
In the coordinates $(a,\ell)$ we have $\varphi(a,\ell)=(ma+t,\ \ell+t)$ with indices taken modulo $135$
and $2$ respectively.

\emph{Case $t$ odd and $m^2\equiv1\pmod{135}$.}  Then $m^2\equiv1\pmod{270}$ and the maps
$x\mapsto x+2j$ and $x\mapsto mx+s$ ($s$ odd) form a group of order $270$.  A map $x\mapsto mx+s$ with
$s$ odd has no fixed point, since $(m-1)x$ is even.  So this group acts regularly, contradicting
Proposition~\ref{prop:regular}.

\emph{Otherwise} some power of $\varphi$ (namely $\varphi$ if $t$ is even, $\varphi^2$ if $t$ is odd)
is level preserving with multiplier $m'\not\equiv1\pmod{135}$.  Composing with a power of $\tau$ we may
assume that $\mu\colon(a,\ell)\mapsto(m'a,\ell)$ is an automorphism, and replacing $m'$ by a power we
may assume that $m'$ has order $2$ or $3$ in $\Z_{135}^{*}\cong\Z_{18}\times\Z_4$.  The elements of
order $2$ are $134,26,109$ and those of order $3$ are $46,91$.  For a block $B$ with
$\mu(B)=B+s$, put $\pi(a,\ell)=(m'a-s,\ell)$, a permutation of the points with $\pi(B)=B$.

\smallskip\noindent\emph{$m'=134=-1$.}  Every block has a level containing at least $3$ points, so it
contains a pure pair.  By Lemma~\ref{lem:transl} with $\mu(Q)-\mu(P)=-(Q-P)$, every block satisfies
$\pi(B)=B$ with $\pi(a)=-a-s$.  This involution has a unique fixed point on each level.  A
$\pi$-invariant set of at least $3$ points on one level contains either the fixed point $y$ and a pair
$\{x,\pi x\}$, giving the repeated difference $x-y=y-\pi x$, or two pairs $\{x,\pi x\},\{z,\pi z\}$,
giving $z-x=\pi x-\pi z$.  Both contradict Lemma~\ref{lem:diff}.

\smallskip\noindent\emph{$m'=26$} ($\equiv-1 \bmod 27$, $\equiv1\bmod 5$).  For $d\in 5\Z_{135}\setminus\{0\}$ we
have $26d=-d$, and these $26$ elements form $13$ pure $\pm$-classes on each level.  By
Lemma~\ref{lem:transl} each such class is covered by a block with $\pi(B)=B$, where
$\pi(a)=26a-s$.

Since $\pi^2(a)=a-27s$ and $B$ is $\pi^2$-invariant, a nonzero translation by $27s\in H$ would force
$5$ points of $B$ into one group; hence $5\mid s$.  Then the fixed points of $\pi$ on a level form one
coset of $H$, i.e.\ they lie in one group, so $B$ contains exactly one fixed point and two
$\pi$-orbits $\{x,\pi x\}$ of length $2$ (each on one level).  Each such orbit contributes one of the
special pure classes ($\pi x-x=25x-s\in5\Z$).  No other pair of $B$ does: if $\{P,Q\}\subseteq B$ is a
pure pair that is not a $\pi$-orbit and $Q-P\in5\Z$, then $\{\pi P,\pi Q\}\subseteq B$ is a different pair
with $\pi Q-\pi P=26(Q-P)=-(Q-P)$, so the same class would be covered twice.  Hence each such block covers exactly $2$ of the $26$ special classes, giving $13$
such orbits.  Every remaining orbit $B$ satisfies $\mu(B)\neq B+s$, and these orbits are permuted by
$\mu$ in pairs.  But $26-13=13$ is odd, a contradiction.

\smallskip\noindent\emph{$m'=109$} ($\equiv1 \bmod 27$, $\equiv-1\bmod 5$).  Now $109d=d$ for
$d\in5\Z_{135}$, and these differences give $13+13$ pure classes and $26$ mixed differences.  By
Lemma~\ref{lem:transl} each is covered by a block with $\pi(B)=B$, $\pi(a)=109a-s$.

Here $\pi^2$ is the translation by $-110s$, of order dividing $27$; invariance of a $5$-set forces
$27\mid s$.  Then a non-fixed point $x$ and $\pi x$ differ by $108x-s\in 27\Z=H$, so they lie in the same
group.  Hence $B$ consists of fixed points of $\pi$, all congruent modulo $5$, and all $10$ differences
of $B$ lie in $5\Z_{135}$.  These blocks therefore cover the $13+13+26=52$ special differences in
blocks of $10$, which is impossible.

\smallskip\noindent\emph{$m'\in\{46,91\}$} (order $3$, $\equiv1\bmod 9$).  Here $m'd=d$ for all
$d\in3\Z_{135}$.  Among any $5$ points two have $a$-coordinates congruent modulo $3$, so by
Lemma~\ref{lem:transl} every block satisfies $\pi(B)=B$ with $\pi(a)=m'a-s$.

As $\pi^3$ is the translation by $-3s$, we get $45\mid s$.  The fixed points of $\pi$ then form a
residue class modulo $3$.  The other $\pi$-orbits are $3$-cycles $\{x,x+c,x+2c\}$ with $c\in45\Z$,
which contain a repeated difference.  Hence every block lies in a residue class modulo $3$, and the
differences not divisible by $3$ are never covered, a contradiction.
\end{proof}

\section{\texorpdfstring{No affine symmetry over $\F_{27}$}{No affine symmetry over F27}}\label{sec:F27}

The field $\F_{27}$ suggests much smaller symmetric searches.  Let $S$ be the group of nonzero squares
of $\F_{27}$ ($|S|=13$; note $-1\notin S$ since $27\equiv3\pmod 4$).

\begin{proposition}\label{prop:agl}
Let $X$ be a set of size $10$.  There is no $5$-GDD of type $10^{27}$ on $\F_{27}\times X$ with groups
$\{u\}\times X$ that is invariant under the group
$\Gamma=\{(u,a)\mapsto(qu+t,a): q\in S,\ t\in\F_{27}\}$ of order $351$.
\end{proposition}

\begin{proof}
A nontrivial element of $\Gamma$ permutes $\F_{27}$ with cycle type $3^9$ or $1^1 13^2$, so it cannot
stabilise a block (whose five points have distinct first coordinates).  Hence all block orbits are
full and there are $3510/351=10$ base blocks.

For points $P=(u,a)$, $Q=(w,c)$ with $u\neq w$, exactly one of $w-u$ and $u-w$ is a square.  Order the
pair so that $w-u\in S$ and call $(a,c)\in X^2$ its \emph{label}.  $\Gamma$ preserves labels and acts
regularly on ordered pairs $(u,w)$ with $w-u\in S$.  So the $100$ orbits of cross pairs correspond to
the $100$ labels, and the $100$ pairs inside the $10$ base blocks must carry every label in $X^2$
exactly once.

Fix $a\in X$ and let $n_a$ be the number of base-block points with second coordinate $a$.  Each such
point lies in $4$ pairs, and $a$ occurs $20$ times as a coordinate of a label (the diagonal label
$(a,a)$ counted twice), so $n_a=5$.  The label $(a,a)$ occurs once, so exactly one base block contains
$a$ twice and no block contains it three times.  A block containing two values twice would contain $4$
pairs with labels in $\{(a,b),(b,a)\}$, but only $2$ are allowed.  Hence each of the $10$ base blocks
has exactly one repeated value, and we may write the blocks as $B_a$ ($a\in X$) with values
$\{a,a\}\cup S_a$, $|S_a|=3$, $a\notin S_a$.

Draw an arc $a\to b$ for $b\in S_a$.  The labels $(a,b),(b,a)$ ($a\neq b$) together occur exactly twice:
$B_a$ contributes $2$ if $b\in S_a$, $B_b$ contributes $2$ if $a\in S_b$, and every $B_c$ with
$a,b\in S_c$ contributes $1$.  Hence there are no $2$-cycles.  Every vertex has out-degree $3$ and, since
$n_a=5$, in-degree $3$.  The complement $\bar G$ of the underlying graph is therefore a cubic graph on
$10$ vertices.  Each $S_c$ spans a triangle of $\bar G$, and every edge of $\bar G$ lies in exactly two
of the ten triangles $S_c$ (counted with multiplicity).

If a triangle $T=S_c=S_{c'}$ is repeated, a vertex $u\in T$ has a third edge $ux$; a triangle through
$ux$ must use a second edge at $u$, which is an edge of $T$ and is already covered twice.  So the ten
triangles are distinct.  Let $uv$ be an edge lying in the distinct triangles $uvw_1$ and $uvw_2$.  As
$\bar G$ is cubic, $u$ has neighbours $v,w_1,w_2$ and $v$ has neighbours $u,w_1,w_2$.  The edge $uw_1$
lies in two triangles, whose third vertices are common neighbours of $u$ and $w_1$, hence lie in
$\{v,w_2\}$; so $w_1w_2$ is an edge.  Then $\{u,v,w_1,w_2\}$ induces a $K_4$, which is a connected
component because the graph is cubic.  So every component of $\bar G$ is $K_4$, which is impossible on
$10$ vertices.
\end{proof}

\begin{remark}
Similarly there is no design on $\F_{27}\times\Z_{10}$ (groups $\{u\}\times\Z_{10}$, levels $x \bmod 2$)
invariant under $(u,x)\mapsto(u+t,x+2j)$ and $(u,x)\mapsto(\omega^2u,x)$ with $\omega$ primitive.
Such a design would have $2$ base blocks.  The $65$ pure level-$0$ $\pm$-differences split into
$5$ orbits of the multiplier, which would require $\binom{e_1}{2}+\binom{e_2}{2}=5$, and no two
binomial coefficients $\binom{e}{2}$ ($0\le e\le 5$) sum to $5$.
\end{remark}

\section{Computer search}\label{sec:search}

All searches used the framework of Section~\ref{sec:framework} (or its equivalent forms
$\Z_{27}\times\Z_5\times\{0,1\}$ and $\F_{27}\times\Z_5\times\{0,1\}$).  The following methods were tried:
\begin{itemize}
\item simulated annealing and min-conflicts tabu search on the $26$ base blocks, including moves that
  shift the level-$1$ part of a block (which changes only mixed differences) and constraint weighting;
\item a two-phase method: first the level-$0$ and level-$1$ shapes (pure differences, found instantly),
  then pairing, signs and shifts so that the mixed differences tile $\Z_{135}\setminus H$
  (exact cover with $\approx 21\,000$ rows, attacked with dancing links and with the SAT solver CaDiCaL);
\item a projection-and-lift method: first the $\Z_{27}$-projection, then labels in $\Z_5\times\{0,1\}$
  making every difference class bijective (local search, CaDiCaL, OR-Tools CP-SAT);
\item randomised backtracking that always covers the most constrained uncovered difference, with restarts.
\end{itemize}

The best configurations found leave between $2$ and $4$ of the $260$ difference classes of
Lemma~\ref{lem:diff} uncovered (and the same number covered twice).  Typically all pure differences
are correct and the defect lies among the mixed differences.  No design was found.

As calibration, the same backtracking search finds a design of type $10^9$ in this framework
($\Z_{90}$, $x\mapsto x+2j$, $8$ base blocks) within seconds.  An exhaustive run shows that the
framework contains no design of type $10^5$ (as expected, since such a design is a
$\mathrm{TD}(5,10)$).  However, none of the methods above recovered a design of type $10^{15}$, which is
known to exist in this framework~\cite{Forbes2023}, within comparable running times.  The negative
outcome for $10^{27}$ therefore reflects the difficulty of the search and is \emph{not} evidence of
non-existence.

\section{Conclusion}

The existence of a $5$-GDD of type $10^{27}$ remains open.  Propositions~\ref{prop:regular},
\ref{prop:affine} and~\ref{prop:agl} rule out the most natural symmetric constructions and show that
a direct construction of Forbes type ($\Z_{270}$, $x\mapsto x+2j$) needs $26$ base blocks with no
additional affine symmetry.  Proposition~\ref{prop:signed} gives a necessary condition that any such
construction must satisfy.  Any candidate list of base blocks can be checked by developing them under
$x\mapsto x+2j$ ($0\le j<135$) and verifying that the $35100$ cross-group pairs are each covered exactly
once.

\begin{thebibliography}{9}
\bibitem{Forbes2023}
A.~D.~Forbes, Group divisible designs with block size five,
\emph{Australas.\ J.\ Combin.} \textbf{87}(1) (2023), 1--11.
\bibitem{ForbesSupp}
A.~D.~Forbes, Group divisible designs with block size five: direct constructions,
arXiv:2211.14124 (2023).
\bibitem{WeiGe}
H.~Wei and G.~Ge, Some more 5-GDDs, 4-frames and 4-RGDDs,
\emph{Discrete Math.} \textbf{336} (2014), 7--21.
\end{thebibliography}

\end{document}
